\chapter{Các công trình nghiên cứu liên quan}
\label{Chapter2}
\setcounter{secnumdepth}{5}

\section{Phân đoạn ảnh y khoa truyền thống}

Phân đoạn ảnh y khoa là một bài toán quan trọng trong phân tích ảnh y tế, với mục tiêu xác định chính xác vùng cơ quan, tổn thương hoặc bất thường bệnh lý trên ảnh. Trong các hệ thống hỗ trợ chẩn đoán, kết quả phân đoạn có thể giúp bác sĩ định lượng kích thước tổn thương, theo dõi tiến triển bệnh và lập kế hoạch điều trị. Vì vậy, yêu cầu đối với mô hình phân đoạn ảnh y khoa không chỉ là dự đoán đúng vùng quan tâm, mà còn cần giữ được ranh giới không gian ở mức chi tiết.

Các mô hình học sâu ban đầu cho phân đoạn ảnh y khoa chủ yếu dựa trên mạng nơ-ron tích chập (Convolutional Neural Networks -- CNNs). Những kiến trúc như FCN, U-Net và UNet++ đã trở thành nền tảng phổ biến nhờ khả năng trích xuất đặc trưng cục bộ và khôi phục bản đồ phân đoạn ở độ phân giải cao \cite{li_lvit_2023}. Đặc biệt, kiến trúc dạng encoder--decoder với các kết nối tắt giúp mô hình vừa học được đặc trưng ngữ nghĩa sâu, vừa giữ lại thông tin không gian từ các tầng nông.

Tuy nhiên, các phương pháp dựa trên CNN vẫn tồn tại một số hạn chế. Thứ nhất, phạm vi tiếp nhận cục bộ của phép tích chập khiến mô hình khó nắm bắt đầy đủ quan hệ toàn cục trong ảnh, nhất là khi vùng tổn thương có hình dạng phức tạp hoặc xuất hiện rải rác. Thứ hai, hiệu quả của mô hình phụ thuộc mạnh vào số lượng và chất lượng mặt nạ gán nhãn. Trong y khoa, việc tạo nhãn pixel-level thường tốn nhiều thời gian, cần chuyên gia có kinh nghiệm và dễ có sai khác giữa các người gán nhãn \cite{li_lvit_2023}. Do đó, dù các mô hình CNN đạt hiệu quả tốt trong nhiều thiết lập có giám sát, khả năng mở rộng sang dữ liệu mới hoặc môi trường lâm sàng thực tế vẫn bị giới hạn.

Sự xuất hiện của Vision Transformer mở ra hướng tiếp cận mới cho phân đoạn ảnh nhờ khả năng mô hình hóa quan hệ dài hạn giữa các vùng ảnh. Transformer có thể xem ảnh như một chuỗi token và học sự phụ thuộc toàn cục thông qua cơ chế self-attention. Tuy vậy, trong ảnh y khoa, chỉ sử dụng Transformer thuần túy có thể làm mất một phần chi tiết cục bộ quan trọng ở biên tổn thương. Vì vậy, nhiều nghiên cứu gần đây có xu hướng kết hợp CNN và Transformer để tận dụng đồng thời khả năng giữ chi tiết cục bộ của CNN và khả năng mô hình hóa ngữ cảnh toàn cục của Transformer.

\section{Học bán giám sát trong phân đoạn ảnh y khoa}

Một trong những thách thức lớn nhất của phân đoạn ảnh y khoa là sự khan hiếm dữ liệu có nhãn. Khác với ảnh tự nhiên, nơi dữ liệu chú thích có thể được thu thập ở quy mô lớn, dữ liệu y khoa thường bị giới hạn bởi quyền riêng tư, chi phí thu thập và yêu cầu chuyên môn khi gán nhãn. Điều này thúc đẩy các phương pháp học bán giám sát, trong đó mô hình học từ một phần nhỏ dữ liệu có nhãn và một lượng lớn dữ liệu không nhãn.

Các phương pháp bán giám sát thường khai thác ý tưởng rằng dữ liệu không nhãn vẫn chứa thông tin về phân bố ảnh. Một số hướng tiếp cận phổ biến gồm tạo nhãn giả, ràng buộc tính nhất quán giữa các dự đoán, hoặc sử dụng nhiều mô hình/nhánh dự đoán để kiểm tra lẫn nhau. Trong bối cảnh được tổng hợp bởi LViT, các phương pháp như DTC, PLCT và MC-Net+ được xem là những đại diện cho xu hướng tận dụng dữ liệu không nhãn nhằm giảm phụ thuộc vào mặt nạ thủ công \cite{li_lvit_2023}.

Tuy nhiên, học bán giám sát trong ảnh y khoa cũng có rủi ro riêng. Nếu nhãn giả có chất lượng thấp, mô hình có thể học lại các dự đoán sai và làm suy giảm hiệu quả tổng quát hóa. Vấn đề này đặc biệt nghiêm trọng khi vùng tổn thương nhỏ, ranh giới mờ hoặc dữ liệu có sự khác biệt lớn giữa các ca bệnh. Do đó, một quy trình bán giám sát hiệu quả cần kiểm soát chất lượng nhãn giả và cập nhật nhãn giả một cách ổn định trong quá trình huấn luyện.

Trong công trình LViT, Li và cộng sự đề xuất cơ chế Exponential Pseudo label Iteration (EPI) nhằm cải thiện chất lượng nhãn giả trong thiết lập bán giám sát \cite{li_lvit_2023}. Thay vì sử dụng trực tiếp nhãn giả từ một bước dự đoán đơn lẻ, EPI cập nhật nhãn giả theo hướng lặp lại và làm ổn định quá trình học. Cơ chế này giúp giảm tác động của các dự đoán nhiễu, đồng thời hỗ trợ mô hình khai thác dữ liệu không nhãn hiệu quả hơn. Đây là một thành phần quan trọng để LViT đạt kết quả cạnh tranh ngay cả khi chỉ sử dụng một phần dữ liệu có nhãn.

\section{Phân đoạn ảnh theo tham chiếu ngôn ngữ}

Phân đoạn ảnh theo tham chiếu ngôn ngữ (Referring Image Segmentation -- RIS) là bài toán phân đoạn đối tượng hoặc vùng ảnh được mô tả bởi một câu văn bản tự nhiên. Khác với phân đoạn ngữ nghĩa truyền thống, trong đó các lớp cần phân đoạn được định nghĩa cố định, RIS cho phép người dùng chỉ định vùng quan tâm bằng ngôn ngữ. Do đó, cùng một ảnh có thể tạo ra các mặt nạ khác nhau tùy thuộc vào nội dung câu mô tả.

Trên ảnh tự nhiên, nhiều phương pháp RIS ban đầu thường trích xuất đặc trưng ảnh và đặc trưng văn bản bằng hai bộ mã hóa riêng biệt, sau đó dung hợp chúng ở giai đoạn giải mã. Cách tiếp cận này có ưu điểm là dễ thiết kế theo kiến trúc modular, nhưng việc dung hợp muộn có thể làm hạn chế mức độ tương tác sâu giữa token văn bản và đặc trưng thị giác. LAVT được đề xuất để giải quyết hạn chế này bằng cách đưa cơ chế tương tác ngôn ngữ--thị giác vào sớm trong bộ mã hóa Vision Transformer \cite{yang_lavt_2021}. Cụ thể, mô hình sử dụng Pixel-Word Attention Module (PWAM) để cho phép đặc trưng ảnh tương tác với đặc trưng từ ngay trong quá trình mã hóa, từ đó tạo ra biểu diễn đa phương thức giàu ngữ cảnh hơn.

Ý tưởng của LAVT cho thấy tầm quan trọng của việc dung hợp sớm giữa ảnh và văn bản. Khi đặc trưng thị giác được điều hướng bởi câu mô tả ngay từ các tầng mã hóa, mô hình có thể tập trung tốt hơn vào đối tượng được nhắc đến, thay vì phải dựa vào bộ giải mã để kết hợp hai nguồn thông tin ở giai đoạn muộn. Đây là một hướng tiếp cận có ảnh hưởng đến các mô hình phân đoạn có điều kiện theo ngôn ngữ sau này, trong đó bao gồm cả các mô hình ứng dụng cho ảnh y khoa.

\section{Phân đoạn ảnh y khoa có tham chiếu văn bản}

Khi chuyển từ ảnh tự nhiên sang ảnh y khoa, bài toán phân đoạn theo tham chiếu ngôn ngữ trở nên khó hơn. Ảnh y khoa thường có độ tương phản thấp, ranh giới tổn thương không rõ, cấu trúc giải phẫu phức tạp và sự khác biệt lớn giữa các phương thức ảnh như X-ray, CT hoặc MRI. Ngoài ra, mô tả văn bản trong y khoa không chỉ gọi tên đối tượng, mà còn có thể chứa thông tin về vị trí, mức độ, phạm vi và đặc điểm bệnh lý.

LViT là một trong những công trình tiêu biểu đưa thông tin văn bản vào phân đoạn ảnh y khoa \cite{li_lvit_2023}. Mô hình được thiết kế theo hướng kết hợp kiến trúc CNN và Transformer, đồng thời sử dụng văn bản như một nguồn tín hiệu bổ sung để định hướng quá trình phân đoạn. Trong LViT, thông tin ảnh được xử lý qua các nhánh thị giác nhằm giữ đặc trưng không gian, trong khi thông tin văn bản được mã hóa và đưa vào quá trình học biểu diễn. Pixel-Level Attention Module (PLAM) giúp mô hình bảo toàn đặc trưng cục bộ, còn Language-Vision loss hỗ trợ khai thác thông tin văn bản trong thiết lập có dữ liệu không nhãn.

Điểm đáng chú ý của LViT là mô hình không chỉ sử dụng văn bản trong giai đoạn suy luận, mà còn đưa văn bản vào quá trình huấn luyện để cải thiện chất lượng biểu diễn. Với ảnh y khoa, chú thích văn bản có thể cung cấp tín hiệu ngữ nghĩa quan trọng về vùng bệnh lý, đặc biệt khi mặt nạ phân đoạn không đầy đủ hoặc số lượng nhãn bị giới hạn. Vì vậy, LViT có thể được xem là cầu nối giữa phân đoạn ảnh y khoa truyền thống và các mô hình phân đoạn đa phương thức dựa trên ngôn ngữ.

Tuy nhiên, LViT giả định rằng chú thích văn bản đi kèm ảnh là tương đối đầy đủ. Trong thực tế lâm sàng, báo cáo hoặc mô tả y khoa có thể thiếu một số thuộc tính quan trọng, chẳng hạn như vị trí, phạm vi hoặc số lượng tổn thương. ARSeg được đề xuất nhằm xử lý vấn đề này trong bài toán Medical Image Referring Segmentation (MIRS) \cite{Wang_Lin_Yan}. Mô hình xây dựng các đặc trưng đặc thù theo thuộc tính và sử dụng cơ chế tương tác chéo đa phương thức để giảm sự phụ thuộc vào văn bản đầy đủ. Ngoài ra, các hàm mất mát cân bằng thuộc tính được thiết kế để tránh việc mô hình thiên lệch về những thuộc tính xuất hiện thường xuyên hơn.

So với LViT, ARSeg nhấn mạnh tính bền vững của mô hình khi prompt văn bản bị thiếu hoặc mất cân bằng. Đây là một hướng mở rộng quan trọng, bởi dữ liệu y khoa thực tế thường không có cấu trúc hoàn hảo như trong bộ dữ liệu chuẩn. Nhìn chung, LViT tập trung vào việc chứng minh hiệu quả của việc đưa văn bản vào phân đoạn ảnh y khoa, còn ARSeg mở rộng vấn đề sang bối cảnh văn bản không hoàn chỉnh và gần hơn với thực tế lâm sàng.

\section{Phân đoạn linh hoạt theo prompt và thế giới mở}

Bên cạnh các mô hình phân đoạn y khoa có điều kiện theo văn bản, một nhánh nghiên cứu khác hướng đến khả năng phân đoạn linh hoạt trong giai đoạn suy luận. Thay vì huấn luyện mô hình cho một tập lớp cố định, các phương pháp này cho phép mô hình nhận prompt dưới nhiều dạng khác nhau, chẳng hạn như văn bản, ảnh ví dụ, hộp bao hoặc điểm gợi ý.

CLIPSeg là một ví dụ tiêu biểu cho hướng phân đoạn dựa trên prompt \cite{luddecke_image_2021}. Công trình này mở rộng mô hình CLIP bằng một bộ giải mã Transformer để tạo bản đồ phân đoạn từ prompt văn bản hoặc ảnh ví dụ. Nhờ sử dụng không gian biểu diễn chung của CLIP, CLIPSeg có thể thực hiện nhiều tác vụ như referring expression segmentation, zero-shot segmentation và one-shot segmentation trong cùng một mô hình. Điểm mạnh của hướng tiếp cận này là tính linh hoạt: người dùng có thể mô tả vùng cần phân đoạn bằng ngôn ngữ tự nhiên hoặc cung cấp một ảnh tham chiếu thay vì phải huấn luyện lại mô hình cho từng lớp mới.

Grounded SAM tiếp tục mở rộng ý tưởng phân đoạn thế giới mở bằng cách kết hợp nhiều mô hình nền tảng \cite{ren_grounded_2024}. Trong pipeline này, Grounding DINO được sử dụng để phát hiện vùng ảnh tương ứng với văn bản đầu vào, sau đó các hộp giới hạn được dùng làm prompt cho Segment Anything Model (SAM) để sinh mặt nạ phân đoạn. Cách kết hợp này tận dụng thế mạnh riêng của từng mô hình: Grounding DINO có khả năng định vị đối tượng theo từ vựng mở, còn SAM có khả năng tạo mặt nạ chất lượng cao từ prompt trực quan.

Mặc dù CLIPSeg và Grounded SAM chủ yếu được phát triển cho ảnh tự nhiên hoặc bối cảnh tổng quát, các công trình này cho thấy xu hướng quan trọng của lĩnh vực phân đoạn: mô hình ngày càng ít phụ thuộc vào tập nhãn cố định và ngày càng tương tác nhiều hơn với prompt. Đối với ảnh y khoa, xu hướng này gợi mở khả năng xây dựng các hệ thống phân đoạn linh hoạt, trong đó bác sĩ có thể chỉ định vùng quan tâm bằng văn bản hoặc gợi ý trực quan. Tuy nhiên, việc chuyển các mô hình thế giới mở sang y khoa vẫn cần xử lý khác biệt miền dữ liệu, tiêu chuẩn an toàn và độ chính xác lâm sàng.

\section{Học biểu diễn đa phương thức ảnh--văn bản trong y khoa}

Một hướng nghiên cứu liên quan mật thiết đến LViT là học biểu diễn đa phương thức giữa ảnh y khoa và văn bản y khoa. Thay vì chỉ huấn luyện mô hình cho một tác vụ phân đoạn cụ thể, các phương pháp này thường tiền huấn luyện bộ mã hóa ảnh và văn bản trên các cặp ảnh--báo cáo, sau đó chuyển giao biểu diễn đã học sang các tác vụ như phân loại, truy hồi, hỏi đáp, sinh báo cáo hoặc phân đoạn.

MedViLL là một công trình tiêu biểu trong hướng tiền huấn luyện thị giác--ngôn ngữ cho y khoa \cite{moon_multi-modal_2021}. Mô hình sử dụng kiến trúc dựa trên BERT kết hợp cơ chế attention đa phương thức để học từ ảnh X-quang và báo cáo không cấu trúc. Thông qua các mục tiêu tiền huấn luyện đa phương thức, MedViLL hướng đến việc cải thiện hiệu quả trên nhiều tác vụ hạ nguồn như phân loại chẩn đoán, truy hồi ảnh--báo cáo, hỏi đáp y khoa và sinh báo cáo. Công trình này cho thấy văn bản y khoa không chỉ là chú thích phụ trợ, mà còn có thể đóng vai trò quan trọng trong việc hình thành biểu diễn thị giác có ý nghĩa lâm sàng.

MGCA tiếp cận bài toán học biểu diễn đa phương thức ở mức chi tiết hơn bằng cách căn chỉnh ảnh và báo cáo X-quang ở nhiều cấp độ \cite{wang_multi-granularity_2022}. Cụ thể, mô hình xem xét sự tương ứng giữa ảnh và văn bản ở cấp độ bệnh lý, cấp độ instance và cấp độ token/vùng cục bộ. Cách căn chỉnh đa hạt này giúp bộ mã hóa hình ảnh học được cả ngữ nghĩa tổng thể lẫn chi tiết cục bộ liên quan đến bệnh lý. Khi chuyển sang các tác vụ hạ nguồn, đặc biệt là phân đoạn ngữ nghĩa, biểu diễn đã được tiền huấn luyện có thể giúp mô hình đạt hiệu quả tốt hơn ngay cả khi số lượng dữ liệu tinh chỉnh bị hạn chế.

Các công trình như MedViLL và MGCA cung cấp nền tảng quan trọng cho các mô hình phân đoạn đa phương thức. Nếu LViT khai thác văn bản trực tiếp trong kiến trúc phân đoạn, thì các mô hình tiền huấn luyện đa phương thức tập trung vào việc học không gian biểu diễn chung trước khi áp dụng vào tác vụ hạ nguồn. Hai hướng tiếp cận này có thể bổ sung cho nhau: biểu diễn ảnh--văn bản được tiền huấn luyện tốt có thể cải thiện chất lượng đặc trưng đầu vào, trong khi kiến trúc phân đoạn như LViT giúp tận dụng trực tiếp tín hiệu văn bản ở cấp độ pixel.

\section{Tổng kết}

Các công trình liên quan cho thấy phân đoạn ảnh y khoa đang chuyển từ các mô hình có giám sát dựa trên ảnh đơn thuần sang các mô hình linh hoạt hơn, có khả năng tận dụng dữ liệu không nhãn và thông tin ngôn ngữ. Các mô hình CNN như U-Net đặt nền móng cho phân đoạn ảnh y khoa nhờ khả năng giữ chi tiết không gian, trong khi Transformer giúp mở rộng khả năng mô hình hóa quan hệ toàn cục. Học bán giám sát giải quyết phần nào vấn đề thiếu dữ liệu gán nhãn, nhưng cần kiểm soát chất lượng nhãn giả để tránh lan truyền sai số.

Trong bối cảnh phân đoạn theo văn bản, LAVT chứng minh hiệu quả của dung hợp sớm giữa ảnh và ngôn ngữ trên ảnh tự nhiên, còn LViT đưa ý tưởng này vào ảnh y khoa bằng cách kết hợp kiến trúc CNN--Transformer, chú thích văn bản, PLAM, Language-Vision loss và EPI \cite{yang_lavt_2021,li_lvit_2023}. Các hướng mở rộng như ARSeg cho thấy cần quan tâm đến tính không đầy đủ của prompt trong thực tế lâm sàng \cite{Wang_Lin_Yan}. Song song với đó, CLIPSeg, Grounded SAM, MedViLL và MGCA cho thấy tiềm năng của prompt, mô hình nền tảng và học biểu diễn đa phương thức trong việc xây dựng các hệ thống phân đoạn linh hoạt hơn \cite{luddecke_image_2021,ren_grounded_2024,moon_multi-modal_2021,wang_multi-granularity_2022}.

Từ các phân tích trên, có thể thấy LViT là một công trình phù hợp với xu hướng kết hợp thị giác và ngôn ngữ cho phân đoạn ảnh y khoa. Mô hình vừa kế thừa ưu điểm của các kiến trúc phân đoạn truyền thống, vừa khai thác văn bản và dữ liệu không nhãn để cải thiện hiệu quả trong điều kiện dữ liệu y khoa hạn chế. Đây cũng là lý do báo cáo lựa chọn LViT làm đối tượng chính để phân tích, tái hiện và đánh giá trong các chương tiếp theo.